\documentclass[../DoAn.tex]{subfiles}
\begin{document}

\begin{center}
    \Large{\textbf{ABSTRACT}}\\
\end{center}
\vspace{1cm}
Video lectures have become a popular form of online learning thanks to their low cost and broad accessibility. However, because the format is inherently one-directional, it tends to leave learners passive and prone to misjudging how well they have understood the material, while comprehension checks are usually decoupled from the act of watching. Existing platforms either fail to effectively support interactive activities within the lecture itself, or still require instructors to design each activity by hand, which is labor-intensive.

To address these limitations, the project adopts an approach that leverages artificial intelligence, in particular large language models capable of understanding the content and proposing questions closely aligned with the lecture, to automatically generate interactive content directly from the video. This approach reduces the authoring burden on instructors while increasing learners' active engagement.

On this basis, the project develops Dyadia, a video-based learning platform integrated with artificial intelligence. The platform lets instructors organize content into courses and lessons, upload videos, automatically generate interactive activities, and manage learners by role. For learners, each lecture becomes an interactive session with automatic grading and immediate feedback during viewing, while instructors can monitor the learning outcomes of their class.

The main contribution of this project is the proposal and implementation of a workflow that converts passive lecture videos into active lessons at a low authoring cost. The system has been built with a complete end-to-end flow, spanning video upload, content analysis and interactive-activity generation, through to learning, automatic grading, feedback, and progress tracking.

\end{document}
